\chapter{Yapay Zeka Destekli Hata Analizi}

\section{AiFailureAnalyzer Mimarisi}
Test fail olduğunda \texttt{AiFailureAnalyzer.analyze()} çağrılır. İki katmanlı analiz uygulanır:
\begin{enumerate}
  \item Flask ML servisi: flaky/stabil tahmin (\texttt{/predict} endpoint).
  \item Groq LLM: Türkçe kök neden, önerilen aksiyon ve risk değerlendirmesi.
\end{enumerate}

\section{Groq LLM ile Kök Neden Analizi}
Groq OpenAI uyumlu API üzerinden \texttt{llama-3.1-8b-instant} modeli kullanılır. Analiz çıktısı şu başlıkları içerir:
\begin{itemize}
  \item Kök Neden
  \item URL / Assertion Detayları
  \item Önerilen Aksiyon
  \item Risk Değerlendirmesi
\end{itemize}

\section{Flask ML Modeli (Flaky/Stabil Tahmin)}
\texttt{ai-analysis/app.py} servisi geçmiş test metadata'sından eğitilmiş \texttt{qa\_model.pkl} modelini sunar. Modül, teknik, standart ve süre bilgileri feature olarak kullanılır.

\section{FailureLogWriter — Teknik Hata Logları}
Fail anında yalnızca teknik bilgi içeren log dosyası \texttt{target/test-history/logs/} altına yazılır. Stack trace, assertion mesajı ve test metadata'sı içerir.

\section{Türkçe Analiz Çıktı Formatı}
Dashboard Hata Analizi ve Raporlanan Hatalar bölümlerinde AI notu Türkçe gösterilir. \texttt{[AI-SKIPPED]} etiketi API anahtarı yoksa veya servis kapalıysa kullanılır.

% TODO: Gerçek FPW-EG-01 fail analizi örneğini listing olarak ekle.
